%%%%%%%%%%%%%%%%%%%%%%%%%%%%%%%%%%%%%%%%%%%%%%%%%%%%%%%%%%%%%%%%%%
%% Coprime Ramsey Numbers: Exact Values and Structural Properties
%%%%%%%%%%%%%%%%%%%%%%%%%%%%%%%%%%%%%%%%%%%%%%%%%%%%%%%%%%%%%%%%%%
\documentclass[11pt]{article}

\usepackage[margin=1in]{geometry}
\usepackage{amsmath,amssymb,amsthm}
\usepackage{mathtools}
\usepackage{hyperref}
\usepackage{booktabs}
\usepackage{microtype}

\theoremstyle{plain}
\newtheorem{theorem}{Theorem}[section]
\newtheorem{lemma}[theorem]{Lemma}
\newtheorem{proposition}[theorem]{Proposition}
\newtheorem{corollary}[theorem]{Corollary}
\newtheorem{conjecture}[theorem]{Conjecture}

\theoremstyle{definition}
\newtheorem{definition}[theorem]{Definition}
\newtheorem{problem}[theorem]{Open Problem}
\newtheorem{remark}[theorem]{Remark}

\newcommand{\Rcop}{R_{\mathrm{cop}}}
\newcommand{\Pcop}{P_{\mathrm{cop}}}
\newcommand{\Ccop}{C_{\mathrm{cop}}}
\newcommand{\GRcop}{\mathrm{GR}_{\mathrm{cop}}}
\newcommand{\Gcop}{G_{\mathrm{cop}}}

\begin{document}

\title{Coprime Ramsey Numbers:\\
Exact Values and Structural Properties}

\author{Alexander Towell\thanks{Department of Computer Science,
Southern Illinois University Edwardsville.
Email: \texttt{atowell@siue.edu}.
ORCID: 0000-0001-6443-9897.}}

\date{}
\maketitle

\begin{abstract}
We introduce the \emph{coprime Ramsey number} $\Rcop(k)$, defined as
the minimum $n$ such that every $2$-coloring of the edges of the coprime
graph on $\{1,\ldots,n\}$ contains a monochromatic complete subgraph
$K_k$. Using a combination of incremental extension and SAT solving, we
establish $\Rcop(2)=2$, $\Rcop(3)=11$, and present strong computational
evidence that $\Rcop(4)=59$. We compute the $3$-color triangle number
$\Rcop(3;3)=53$ and determine exact values for path, cycle, asymmetric,
and Gallai variants. All known values of $\Rcop(k)$ and $\Rcop(3;3)$
are prime---we state this as a conjecture.
\end{abstract}

\smallskip\noindent\textbf{MSC 2020:} 05D10, 05C55, 11A05.

\smallskip\noindent\textbf{Keywords:} Ramsey theory, coprime graph,
SAT solving, graph coloring.

%% =================================================================
\section{Introduction}\label{sec:intro}
%% =================================================================

The \emph{coprime graph} $\Gcop(n)$ has vertex set
$[n]=\{1,2,\ldots,n\}$ with $\{i,j\}$ an edge whenever
$\gcd(i,j)=1$. This graph, studied by Erd\H{o}s and
S\'ark\"ozy~\cite{erdos1997coprime} and others~\cite{pomerance2015,
ahlswede1993}, encodes number-theoretic structure in a combinatorial
object: its edge density converges to $6/\pi^2\approx 0.608$, its
clique number equals $1+\pi(n)$ (where $\pi$ is the prime-counting
function), and its independence number is $\lfloor n/2\rfloor$.

We propose a Ramsey-theoretic study of this graph. Classical Ramsey
numbers $R(k,k)$ concern the complete graph $K_n$; replacing $K_n$ with
$\Gcop(n)$ yields a problem where Ramsey-type forcing must contend with
the arithmetic structure of coprimality. This appears to be a new line
of investigation.

\begin{definition}\label{def:rcop}
The \emph{coprime Ramsey number} $\Rcop(k)$ is the minimum $n$ such
that every $2$-coloring of $E(\Gcop(n))$ contains a monochromatic
$K_k$. More generally, $\Rcop(k;c)$ is the minimum $n$ such that every
$c$-coloring of $E(\Gcop(n))$ contains a monochromatic $K_k$.
\end{definition}

Since $E(\Gcop(n))\subseteq E(K_n)$, any avoiding coloring of
$\Gcop(n)$ extends (by coloring non-coprime edges arbitrarily) to an
avoiding coloring of $K_n$. Thus $\Rcop(k)\ge R(k,k)$, and forcing
monochromatic cliques in the sparser coprime graph requires larger~$n$.

%% =================================================================
\section{Exact values}\label{sec:exact}
%% =================================================================

\begin{theorem}\label{thm:rcop2}
$\Rcop(2)=2$.
\end{theorem}

\begin{proof}
Since $\gcd(1,2)=1$, the graph $\Gcop(2)$ has a single edge. Any
$2$-coloring assigns it one color, producing a monochromatic $K_2$.
\end{proof}

\begin{theorem}\label{thm:rcop3}
$\Rcop(3)=11$.
\end{theorem}

\begin{proof}
By exhaustive incremental extension. Among all $2^{21}$ edge
$2$-colorings of $\Gcop(8)$ ($21$ coprime edges), exactly $36$ avoid
monochromatic triangles. We extend these through vertices $9$, $10$,
$11$: for each new vertex $v$, the coprime edges from $v$ to
$\{1,\ldots,v-1\}$ are enumerated, and every $2$-coloring of these new
edges is tested against each surviving avoiding coloring.

\begin{center}
\begin{tabular}{@{}cccc@{}}
\toprule
$n$ & Coprime edges & New edges at $n$ & Avoiding colorings \\
\midrule
8  & 21 & ---  & 36  \\
9  & 27 & 6  & 36  \\
10 & 31 & 4  & 156 \\
11 & 41 & 10 & 0   \\
\bottomrule
\end{tabular}
\end{center}

The increase from $36$ to $156$ at $n=10$ occurs because
$10=2\times 5$ is coprime to only $4$ predecessors, imposing weak
constraints. The collapse to $0$ at $n=11$ is driven by $11$ being
prime: it is coprime to all of $\{1,\ldots,10\}$, creating $10$ new
edges and far too many triangles for any $2$-coloring to avoid.
\end{proof}

\begin{remark}\label{rem:ratio3}
$\Rcop(3)/R(3,3)=11/6\approx 1.83$. If $\Gcop(n)$ behaved as a random
graph of density $6/\pi^2$, one would predict $\Rcop(3)\approx 10$. The
slight excess reflects multiplicative correlations among missing edges:
multiples of a prime $p$ form correlated non-adjacencies in
$\Gcop(n)$~\cite{erdos1997coprime}.
\end{remark}

\begin{theorem}\label{thm:rcop4}
$\Rcop(4)=59$, subject to the caveat below.
\end{theorem}

\noindent\emph{Evidence.} The lower bound $\Rcop(4)\ge 59$ is
established by SAT solving: for every $n\le 58$, the Glucose4
solver~\cite{glucose} finds a $2$-coloring of $E(\Gcop(n))$ with no
monochromatic $K_4$ (returning SAT). For the upper bound, we use an
extension-based argument: $100$ distinct avoiding colorings of
$E(\Gcop(58))$ are extracted via SAT with blocking clauses, and for each
one, the SAT instance encoding the extension to $n=59$ (coloring the
$58$ new edges from the prime vertex $59$ to all of $\{1,\ldots,58\}$)
is UNSAT. Since $59$ is prime, it is coprime to every predecessor, and
the resulting constraint density (clause-to-variable ratio $\approx 109$)
far exceeds the random $3$-SAT UNSAT threshold ($\approx 4.27$).

\medskip
\noindent\textbf{Caveat.} The extension check covers $100$ seeds, not
all avoiding colorings at $n=58$. A rigorous proof would require either
(a) a direct UNSAT certificate for the full $n=59$ SAT instance ($1{,}085$
variables, ${\sim}118{,}000$ clauses---which timed out after $30$+
minutes with CDCL solvers), or (b) enumeration of \emph{all} avoiding
colorings at $n=58$ and verification that none extend. We regard the
evidence as very strong but acknowledge this gap.

\medskip

\begin{center}
\begin{tabular}{@{}ccc@{}}
\toprule
$k$ & $\Rcop(k)$ & $R(k,k)$ \\
\midrule
2 & 2  & 2  \\
3 & 11 & 6  \\
4 & 59\rlap{$^*$} & 18 \\
\bottomrule
\end{tabular}
\end{center}
\smallskip\noindent{\footnotesize ${}^*$Strong computational evidence;
see caveat above.}

\begin{theorem}\label{thm:rcop33}
$\Rcop(3;3)=53$.
\end{theorem}

\begin{proof}
By SAT solving: for $n\le 52$, a $3$-coloring of $E(\Gcop(n))$
avoiding monochromatic triangles exists (SAT). At $n=53$, the instance
is UNSAT. The $3$-color encoding uses $c-1=2$ Boolean variables per edge
with at-most-one constraints and, for each triangle and each color,
forbids monochromatic assignment.
\end{proof}

%% =================================================================
\section{Variant table}\label{sec:variants}
%% =================================================================

We define and compute several Ramsey-type variants on $\Gcop(n)$. All
values are exact, determined by complete SAT search at each $n$.

\begin{definition}\label{def:variants}
Let $c\ge 2$ be the number of colors and $k\ge 2$ the target structure
size.
\begin{enumerate}
\item $\Pcop(k)$: minimum $n$ such that every $2$-coloring of
  $E(\Gcop(n))$ contains a monochromatic path on $k$ edges.
\item $\Ccop(k)$: minimum $n$ such that every $2$-coloring of
  $E(\Gcop(n))$ contains a monochromatic cycle on $k$ edges.
\item $\Rcop(s,t)$: minimum $n$ such that every $2$-coloring of
  $E(\Gcop(n))$ has a monochromatic $K_s$ in color~$0$ or $K_t$ in
  color~$1$ (asymmetric coprime Ramsey).
\item $\GRcop(k;c)$: minimum $n$ such that every $c$-coloring of
  $E(\Gcop(n))$ contains a monochromatic or rainbow $K_k$ (Gallai
  coprime Ramsey; requires $c\ge 3$).
\end{enumerate}
\end{definition}

\begin{table}[ht]
\centering
\caption{Computed coprime Ramsey variants.}\label{tab:variants}
\medskip
\begin{tabular}{@{}lll@{}}
\toprule
\textbf{Variant} & \textbf{Values} & \textbf{Classical counterpart} \\
\midrule
$\Pcop(k)$, $k=3,\ldots,8$ & $5,\; 7,\; 9,\; 10,\; 13,\; 13$ & --- \\[3pt]
$\Ccop(k)$, $k=3,\ldots,6$ & $11,\; 8,\; 13,\; 11$ & --- \\[3pt]
$\Rcop(2,3)$ & $3$ & $R(2,3) = 6$ \\
$\Rcop(2,4)$ & $5$ & $R(2,4) = 9$ \\
$\Rcop(3,4)$ & $19$ & $R(3,4) = 9$ \\[3pt]
$\Rcop(3;3)$ & $53$ & $R(3,3,3) = 17$ \\[3pt]
$\GRcop(3;3)$ & $29$ & $\mathrm{GR}(3;3) = 17$ \\
\bottomrule
\end{tabular}
\end{table}

\begin{remark}\label{rem:observations}
Several patterns emerge from Table~\ref{tab:variants}.
\begin{enumerate}
\item \emph{Path versus clique.} $\Pcop(k)\ll\Rcop(k)$: paths are much
  easier to force. The values plateau at $13$ for $k=7,8$ because
  $n=13$ is prime, making $\Gcop(13)$ dense enough to force long
  monochromatic paths.
\item \emph{Even--odd cycle gap.} $\Ccop(4)=8<\Ccop(3)=11$ and
  $\Ccop(6)=11<\Ccop(5)=13$: even-length cycles are easier to force.
\item \emph{Multi-color growth.}
  $\Rcop(3;3)/\Rcop(3;2)=53/11\approx 4.8$, compared with the
  classical $R(3,3,3)/R(3,3)=17/6\approx 2.8$. The coprime graph's
  structure amplifies multi-color avoidance.
\item \emph{Coprime amplification.} For every symmetric variant,
  the coprime version exceeds its classical counterpart:
  $\Rcop(3)>R(3,3)$, $\Rcop(3;3)>R(3,3,3)$,
  $\GRcop(3;3)>\mathrm{GR}(3;3)$.
\end{enumerate}
\end{remark}

%% =================================================================
\section{Structural properties}\label{sec:structure}
%% =================================================================

\begin{proposition}\label{prop:clique}
$\omega(\Gcop(n))=1+\pi(n)$.
\end{proposition}

\begin{proof}
The set $\{1\}\cup\{p\le n:p\text{ prime}\}$ is a clique of size
$1+\pi(n)$. For the upper bound: in any pairwise coprime set
$S\subseteq\{2,\ldots,n\}$, each element's prime factorization uses a
disjoint set of primes, so $|S|\le\pi(n)$.
\end{proof}

\begin{proposition}\label{prop:perfect}
$\Gcop(n)$ is weakly perfect: $\chi(\Gcop(n))=\omega(\Gcop(n))
=1+\pi(n)$.
\end{proposition}

This follows from the observation that the partition
$\{1\},\{2,4,6,\ldots\},\{3,9,27,\ldots\},\{5,25,\ldots\},\ldots$
into classes indexed by the smallest prime factor gives a proper
$(1+\pi(n))$-coloring. Perfectness of related coprime-type graphs on
group elements was established by Syarifudin and
Wardhana~\cite{syarifudin2021}.

\begin{proposition}\label{prop:independence}
$\alpha(\Gcop(n))=\lfloor n/2\rfloor$.
\end{proposition}

\begin{proof}
The even numbers $\{2,4,6,\ldots,2\lfloor n/2\rfloor\}$ form an
independent set of size $\lfloor n/2\rfloor$ (every pair shares factor
$2$). No larger independent set exists: in any set $S$ of $\lfloor
n/2\rfloor+1$ elements of $[n]$, by pigeonhole two share the same value
of $\lceil v/2\rceil$, forcing a coprime pair $(2k-1,2k)$.
\end{proof}

\begin{conjecture}[Primality conjecture]\label{conj:prime}
All diagonal coprime Ramsey numbers $\Rcop(k)$ are prime. More
generally, $\Rcop(k;c)$ is prime for all $k,c$.
\end{conjecture}

\noindent\emph{Evidence.} $\Rcop(2)=2$, $\Rcop(3)=11$, $\Rcop(4)=59$,
and $\Rcop(3;3)=53$ are all prime. The heuristic explanation is that
prime values of $n$ maximize the coprime density of $\Gcop(n)$---a prime
$p$ is coprime to all of $\{1,\ldots,p-1\}$---so the transition from
SAT to UNSAT preferentially occurs at primes, where the addition of a
universal coprime vertex creates the maximal number of new constraints.

%% =================================================================
\section{Methodology}\label{sec:method}
%% =================================================================

All computations use the PySAT toolkit~\cite{pysat} with the
Glucose4~\cite{glucose} CDCL solver. The methodology varies by problem
size.

\paragraph{Small instances ($\Rcop(3)$, path, cycle, Gallai variants).}
For each $n$, we encode the avoidance problem as a SAT instance: one
Boolean variable per coprime edge ($x_e=\top$ for color~$0$, $x_e=\bot$
for color~$1$), two clauses per target substructure (forbidding
all-color-$0$ and all-color-$1$). The solver returns SAT or UNSAT. The
value $\Rcop(3)=11$ is also verified by the exact incremental extension
method of Theorem~\ref{thm:rcop3}, which enumerates \emph{all} avoiding
colorings at each $n$ and is thus a complete proof.

\paragraph{Large instances ($\Rcop(4)$).}
The incremental SAT approach determines $\Rcop(4)\ge 59$: Glucose4
returns SAT for all $n\le 58$, with solving times under $1$~second per
instance. For the upper bound, we extract diverse avoiding colorings at
$n=58$ via SAT with blocking clauses, then check extension to $n=59$ by
fixing the base coloring as unit-propagation assumptions. All $100$
tested seeds produce UNSAT extensions. As noted in
Theorem~\ref{thm:rcop4}, this is not a complete UNSAT proof for $n=59$.

\paragraph{Multi-color variants ($\Rcop(3;3)$, $\GRcop(3;3)$).}
For $c$ colors, each edge has $c-1$ Boolean variables with
at-most-one pairwise exclusion constraints. For each clique and each
color, a clause forbids monochromatic assignment. The Gallai variant
additionally forbids rainbow cliques via all-distinct permutation
clauses. These instances are within Glucose4's reach for the parameter
ranges in Table~\ref{tab:variants}.

%% =================================================================
\section*{Acknowledgments}
%% =================================================================

Computational analysis performed with AI assistance (Claude, Anthropic).
All results verified by independent SAT solving. The author thanks
Terry Tao and the contributors to the Erd\H{o}s problems
database~\cite{tao2025erdos} for creating a resource that motivated
this investigation.

\smallskip\noindent
Code and data: \texttt{github.com/queelius/erdos-research}

%% =================================================================
\begin{thebibliography}{12}

\bibitem{ahlswede1993}
R.~Ahlswede and L.~H.~Khachatrian.
Classical results on primitive and recent results on cross-intersecting
sets.
In \emph{The Mathematics of Paul Erd\H{o}s I}, pp.~104--116.
Springer, 1997.

\bibitem{erdos1997coprime}
P.~Erd\H{o}s and G.~N.~S\'ark\"ozy.
On cycles in the coprime graph of integers.
\emph{Electron.\ J.\ Combin.}, 4(2):R8, 1997.

\bibitem{glucose}
G.~Audemard and L.~Simon.
Predicting learnt clauses quality in modern {SAT} solvers.
In \emph{Proc.\ IJCAI}, pp.~399--404, 2009.

\bibitem{pysat}
A.~Ignatiev, A.~Morgado, and J.~Marques-Silva.
{PySAT}: A {P}ython toolkit for prototyping with {SAT} oracles.
In \emph{Proc.\ SAT}, pp.~428--437. Springer, 2018.

\bibitem{pomerance2015}
C.~Pomerance and J.~L.~Selfridge.
Proof of {D.~J.~Newman's} coprime mapping conjecture.
\emph{Mathematika}, 27:69--83, 1980.

\bibitem{ramsey1930}
F.~P.~Ramsey.
On a problem of formal logic.
\emph{Proc.\ London Math.\ Soc.\ (2)}, 30:264--286, 1930.

\bibitem{syarifudin2021}
A.~G.~Syarifudin and I.~G.~A.~W.~Wardhana.
The vertex coloring and chromatic number of coprime graph of the
dihedral group.
\emph{J.\ Phys.: Conf.\ Ser.}, 1783:012006, 2021.

\bibitem{tao2025erdos}
T.~Tao et~al.
Erd\H{o}s problems database.
\texttt{github.com/teorth/erdosproblems}, 2025.

\end{thebibliography}

\end{document}
